\section{选题背景与意义}

\subsection{选题背景}

青年是整个社会力量中最积极、最有生气的力量\cite{Mao1991}，大学生的价值取向关乎个体成长方向，决定着未来整个社会的价值追求与精神面貌\cite{Xi2014}。步入高等教育阶段的青年正处于价值观确立与定型的关键期，伴随生活空间拓展与社会接触面延伸，他们在个人生活、人际交往、学习劳动与择业以及社会公共参与中不断深化对个人与社会关系的认识，逐步确立是非评判与行为抉择的基本尺度\cite{CPCC2017_Plan}。引导大学生扣好人生的“第一粒扣子”、自觉树立和践行社会主义核心价值观，构成高校落实立德树人根本任务的核心环节\cite{CPCC2017_Ipe}。围绕培养德智体美劳全面发展的社会主义建设者和接班人这一总体目标\cite{Xi2018_Speech}，党的二十届三中全会对完善立德树人机制、健全培育和践行社会主义核心价值观制度机制作出系统部署\cite{CPCC2024_Jueding}，要求新时代高校价值观教育贴合青年的现实生活境遇。按照马克思主义唯物史观的基本观点，青年的现实交往状况由其所处的具体社会关系构成；人的本质在其现实性上是一切社会关系的总和，受一定生产方式制约的现实交往形式直接规定着个体的精神过程与价值观念的生成，并作为现实的社会关系中介对社会意识产生能动反作用\cite{MarxEngels2009}。大学生对主流价值的认知、检验与认同，始终离不开日常人际交往这一传递社会规范、交换生活经验与校准自我定位的现实中介。随着现代生产方式变迁、社会分工持续细化以及资本逻辑向日常生活的渗透，当代大学生的互动形态正突破传统熟人社会的差序格局，借助日益便捷的现代通信与网络媒介，围绕细分兴趣、具体任务与阶段性诉求，同不同专业、院校及社会背景的个体建立起广泛联系。在家庭、行政班级与宿舍等稳定的强关系共同体之外，同跨校学习伙伴、临时项目成员、趣缘社群成员及各类校园搭子开展的轻社交日趋普遍。这类交往区别于开放网络空间中的陌生人单向浏览或匿名舆论站队，以对象可辨认、有实质回应、有一定延续性为基本特征，在情感卷入、背景了解与持续投入上相对有限，呈现出鲜明的弱连接形态。广泛分布于线上线下场景中的弱连接交往，承载着多元异质的生活经验与评价标准，成为考察当代大学生价值观生成与变迁的重要现实切入点。

弱连接交往本身具有浅层结构形态与低约束互动规则的特征。不同生活背景与发展轨迹的交往对象充当着跨越群体边界的信息桥梁，在临时协作与兴趣交流中将熟人圈层之外的经验引入大学生的认知视野，为其跳出封闭思维定势、培育契约意识与边界意识提供了客观条件。当这种缺乏深层情感连带与完整背景了解的浅层交往在校园日常中持续扩展并逐步替代深度交往时，其内含的浅层化、功能化与低义务规则反作用于大学生衡量得失与处理群己关系的尺度，促使部分学生在个人生活价值观、人际交往价值观、学习劳动与择业价值观以及社会公共价值观上的取向，同唯物史观价值立场、社会主义核心价值观及当代育人要求产生现实偏差。在个人生活领域，弱连接对象间相互展示的片面消费与休闲片段容易刺激攀比心理，低门槛的浅层社交消遣逐渐替代需要持续投入耐心的深度阅读与自我修养，使课余生活趋于表层化。在学习劳动与择业领域，高等教育劳动力市场竞争压力下的升学、录用与绩点等结果信息在弱连接交往中集中呈现，隐去了背后的长期劳动积累，由此诱发的向上比较容易促使部分学生在功利速成与消极退缩之间摇摆。在人际交往领域，商品拜物教向日常交往关系渗透，按需匹配、即时结算、随时断联的互动规则诱导主体物化其人际关系，将防御性的成本收益核算带入人际伦理，削弱宿舍与班级等集体生活所需的道义承诺与信任基础；在异地求学带来的家庭空间分离背景下，这种交往尺度进一步向亲属关系延伸，使部分学生对日常缺乏直接实用往来的弱连接亲戚主动减少联系乃至刻意疏离，削弱家庭亲情联结。在社会公共领域，围绕校园事务或青年议题展开的松散群体互动缺乏共同生活根基与稳定责任约束，容易将复杂的公共事务简化为情绪化的二元对立表态，淡化青年的公共共同体归属感，削弱其对公共规则的自觉遵守以及理性、公正、法治的公民素养。

大学生日常交往形态向弱连接交往的延伸及其伴随的多维价值矛盾，要求高校思想政治教育深入学生的现实关系网络与微观生活世界。《关于新时代加强和改进思想政治工作的意见》明确提出把思想政治工作贯通于教育教学与日常管理全过程，贴近青年实际做好人文关怀与心理疏导\cite{CPCC2021_Yijian}。这一制度要求植根于马克思主义关于交往实践与道德养成的科学论断；列宁指出，青年唯有将学习、教育和训练中的每一步骤同广泛的社会实践及共同体行动紧密结合，才能摆脱狭隘眼界的束缚，养成符合历史发展要求的道德品质\cite{Lenin2017}。立足新时代育人环境，《教育强国建设规划纲要（2024—2035年）》进一步明确要坚持显性教育与隐性教育相统一，拓展实践育人与网络育人空间，统筹推进价值引领、实践体验与环境营造\cite{CPCC2025_Gangyao}。当前高校思想政治工作主要通过院系、行政班级与宿舍等固定建制开展熟人式管理，面对跨越校园围墙、依即时诉求聚散的弱连接交往网络，存在感知滞后与覆盖盲区。本研究立足唯物史观关于交往形式对社会意识能动反作用的基本原理，考察弱连接交往影响大学生日常价值判断的作用机理与现实成因，探索切合青年交往特征的协同育人路径，为高校在流动社交环境中实现主流价值的有效引领提供学理支撑。

\subsection{研究意义}

本研究将大学生价值观的生成与变迁置于日常交往结构转型的具体情境中加以考察，有助于深化对思想政治教育过程中社会关系中介及其反作用机理的学理认识。马克思主义关于现实的人及其社会交往的学说表明，受生产方式变迁与资本逻辑影响的日常交往形式，以实践规则的方式能动地反作用于个体的价值意识。既有学术探讨中，社会学的弱关系网络分析与思想政治教育学的价值观培育多各自独立展开：前者侧重考察弱连接交往在传递异质信息与促成职业流动中的工具性效用，常将行动者的价值观念视为预先给定的静态前提；后者偏重探讨主流意识形态的宏观阐释与话语传播，较少从历史唯物主义与政治经济学批判视角剖析弱连接交往这一日常交往形式影响青年具体价值判断的微观作用路径。本研究结合社会网络结构分析与马克思主义社会关系理论，以对象可辨认、有实质回应、有一定延续性界定弱连接交往边界，确立考察大学生四维价值观变迁的规范标尺，深入剖析竞争结果片面展示、商品拜物教渗透下低义务规则的跨场景延伸、浅层陪伴对深度投入的替代以及松散群体互动引发个体价值认知偏移的内在过程，拓展弱连接理论在青年价值观研究领域的解释边界，为思想政治教育学把握日常社会关系变迁与主流价值认同的内在关联提供学理依据。

剖析弱连接交往引发的价值分歧及其生成逻辑，有助于高校准确把握大学生在流动交往中的认知困惑，提升新时代思想政治教育的针对性。习近平总书记强调，做好高校思想政治工作要因事而化、因时而进、因势而新，密切关注青年在日常学习生活与社会交往中的现实关切\cite{Xi2022}。当代大学生在个人生活价值观、人际交往价值观、学习劳动与择业价值观以及社会公共价值观上的观念波动，大量孕育于趣缘社群往来、跨校学习求职互助、校园搭子社交及弱连接亲属互动等日常经历之中。透过表层社交行为识别学生所采纳的隐性评价标准及其背后的人际物化倾向、亲情疏离困扰与竞争压力，能够为思想政治理论课教学改革与日常学生管理提供具体的经验依据。思想政治理论课作为落实立德树人根本任务的关键课程\cite{Xi2020}，需要将青年在现实交往中遭遇的价值困惑引入课堂教学，在回应学生深层关切的过程中落实《关于深化新时代学校思想政治理论课改革创新的若干意见》的部署要求\cite{CPCC2019_Yijian}。围绕成因机制所构建的引导方案，可直接服务于高校党团组织建设、跨学科科研协作、家校协同育人与社会实践育人的优化设计，推动思想政治教育实现叙事转换与日常融入，将主流价值融入青年易于理解与参与的日常生活实践\cite{Shen2021_narrative}，引导大学生运用唯物史观方法辨析人际物化倾向与片面比较心理，在将个人奋斗融入国家现代化建设的实践中实现自由全面发展。